\documentclass[11pt,a4paper]{article}

% ---------------------------------------------------------
% PACKAGES
% ---------------------------------------------------------
\usepackage[utf8]{inputenc}
\usepackage[T1]{fontenc}
\usepackage{lmodern}
\usepackage{amsmath,amssymb,amsfonts,bm}
\usepackage{geometry}
\usepackage{graphicx}
\usepackage{hyperref}
\usepackage{authblk}
\usepackage{cite}
\usepackage{physics}
\usepackage{mathtools}
\usepackage{tensor}
\usepackage{enumitem}

\geometry{margin=2.4cm}

% ---------------------------------------------------------
% TITLE
% ---------------------------------------------------------
\title{\textbf{Companion XII (v49.1)}\\[0.2cm]
\textbf{Relaxation Scale and Freeze-Out Geometry}\\[0.2cm]
\large Foundations of Relaxation-Driven Cyclic Cosmology (RDCC)}

\author[1]{Michael Lehmann}
\affil[1]{\small Independent Researcher, Relaxation-Driven Cyclic Cosmology (RDCC)\\
\texttt{mi.lehmann@gmx.de}}

\date{September 2026}

\begin{document}
\maketitle

% ==========================================================
% ABSTRACT
% ==========================================================
\begin{abstract}
A characteristic relaxation scale \(k_{\mathrm{rel}}\) appears throughout the
phenomenology of Relaxation-Driven Cyclic Cosmology (RDCC). Previous generations
of the framework introduced this scale as a phenomenological quantity governing
growth suppression, gravitational-wave damping, redshift-space distortions, weak
lensing, and gravitational slip. In the present Companion, the relaxation scale
is reinterpreted as a dynamical consequence of the freeze-out geometry of the
relaxation sector. By connecting the global relaxation dynamics to the horizon
structure of the early Universe, the relaxation scale emerges naturally from the
freeze-out epoch. The resulting framework establishes a theoretical bridge
between the dimension-8 inter-sector operator, the infrared relaxation parameter
\(\alpha_{\mathrm{IR}}\), primordial tensor damping, and the observable
large-scale structure of the Universe.
\end{abstract}

\tableofcontents
\newpage

% ==========================================================
\section{Introduction}
% ==========================================================

Several observable sectors of RDCC contain a common characteristic scale. These
include:
\begin{itemize}
  \item the suppression of small-scale structure,
  \item the scale dependence of the growth rate,
  \item redshift-space distortions,
  \item weak gravitational lensing,
  \item gravitational slip,
  \item primordial tensor damping.
\end{itemize}

Historically these phenomena were described using the relaxation scale
\begin{equation}
  k_{\mathrm{rel}}.
\end{equation}
Although observationally useful, the origin of this scale remained unclear. The
purpose of the present Companion is to demonstrate that the relaxation scale
arises naturally from the freeze-out geometry of the relaxation sector and is
therefore not an independent fundamental parameter.

This Companion occupies a central position within the RDCC architecture. It
provides the missing bridge between
\begin{itemize}
  \item the relaxation dynamics developed in Companion VI,
  \item the tensor sector developed in Companion II,
  \item the background cosmology established in Companion XI,
  \item and the perturbation framework of Companion VII.
\end{itemize}

% ==========================================================
\section{Global Relaxation Dynamics}
% ==========================================================

The leading interaction between the visible and mirror sectors is encoded in the
dimension-8 operator
\begin{equation}
  \mathcal{O}_{\mathrm{int}}
  =
  \frac{1}{M^4}
  T_{\mu\nu}^{(+)} T^{(-)\mu\nu}.
\end{equation}
This operator generates the effective relaxation sector of RDCC and gives rise
to the global relaxation parameter
\begin{equation}
  \alpha(t).
\end{equation}

The infrared evolution of the relaxation parameter is described by
\begin{equation}
  \dot{\alpha}
  =
  -
  \Gamma_{\mathrm{IR}}
  \alpha.
\end{equation}
The corresponding solution is
\begin{equation}
  \alpha(t)
  =
  \alpha_{\mathrm{IR}}
  +
  \left(
    \alpha_0 - \alpha_{\mathrm{IR}}
  \right)
  e^{-\Gamma_{\mathrm{IR}} t}.
\end{equation}

The infrared fixed point is
\begin{equation}
  \alpha_{\mathrm{IR}}
  =
  0.0176 \pm 0.0021.
\end{equation}
This single quantity governs the entire RDCC prediction vector, including:
\begin{equation}
  n_s,
  \qquad
  \Delta N_{\mathrm{eff}},
  \qquad
  f\sigma_8,
  \qquad
  \Omega_{\mathrm{GW}}.
\end{equation}

The present Companion extends this framework by showing that the characteristic
relaxation scale also follows from the same dynamics.

% ==========================================================
\section{Connection to the Freeze-Out Picture}
% ==========================================================

Earlier generations of RDCC described relaxation through a Hubble-coupled
relaxation rate
\begin{equation}
  \Gamma(a)
  =
  \beta(a)\, H(a).
\end{equation}
The relaxation parameter evolved according to
\begin{equation}
  \frac{d\alpha}{d\ln a}
  =
  \Gamma(a)\, \alpha.
\end{equation}
Although formulated differently, this description captures the same physical
process as the infrared relaxation equation.

The crucial event is the transition between relaxation-dominated evolution and
standard post-relaxation evolution. This transition defines the freeze-out
scale factor
\begin{equation}
  a_{\mathrm{fr}}.
\end{equation}
In the improved running-coupling picture of RDCC, freeze-out occurs when
\begin{equation}
  \beta(a_{\mathrm{fr}})
  =
  1.
\end{equation}
At this moment,
\begin{equation}
  \Gamma(a_{\mathrm{fr}})
  =
  H(a_{\mathrm{fr}}).
\end{equation}
The freeze-out epoch therefore identifies a distinguished moment in the
evolution of the relaxation sector and provides the geometric origin of the
relaxation scale.
% ==========================================================
\section{Freeze-Out Geometry}
% ==========================================================

The freeze-out epoch marks the end of the relaxation-dominated phase of
cosmological evolution. Before freeze-out, the relaxation sector contributes
significantly to the evolution of the metric response. After freeze-out, the
relaxation parameter approaches its infrared fixed point and remains
effectively constant throughout all observable epochs.

The Phase-1 integrated RDCC framework adopts
\begin{equation}
  a_{\mathrm{fr}} \simeq 10^{-4},
\end{equation}
consistent with the freeze-out analysis of the relaxation sector. At this epoch
the condition
\begin{equation}
  \Gamma(a_{\mathrm{fr}}) = H(a_{\mathrm{fr}})
\end{equation}
separates two cosmological regimes:
\begin{enumerate}[label=\arabic*.]
  \item Relaxation-dominated evolution
  \begin{equation}
    \Gamma \gtrsim H,
  \end{equation}
  \item Post-relaxation evolution
  \begin{equation}
    \Gamma \ll H.
  \end{equation}
\end{enumerate}
The freeze-out surface therefore acts as a geometric boundary within the
cosmological evolution of RDCC. The relaxation scale emerges directly from this
boundary.

% ==========================================================
\section{Horizon Definition of the Relaxation Scale}
% ==========================================================

The earliest explicit appearance of the relaxation scale within the RDCC
framework occurs in the tensor sector. Tensor modes that re-enter the horizon
during the relaxation-dominated epoch experience additional damping, whereas
modes that re-enter later evolve almost identically to their standard
cosmological counterparts. This motivates the definition
\begin{equation}
  k_{\mathrm{rel}}^{(H)}
  =
  a_{\mathrm{fr}} H(a_{\mathrm{fr}}).
  \label{eq:krelH}
\end{equation}
The quantity \(k_{\mathrm{rel}}^{(H)}\) corresponds to the comoving horizon
scale at freeze-out.

Modes satisfying
\begin{equation}
  k < k_{\mathrm{rel}}
\end{equation}
re-enter after the relaxation epoch and therefore experience little or no
relaxation-induced suppression. Conversely,
\begin{equation}
  k > k_{\mathrm{rel}}
\end{equation}
identifies modes that re-enter during the relaxation epoch and consequently
acquire additional damping. The relaxation scale therefore separates the
perturbation spectrum into two dynamically distinct sectors.

% ==========================================================
\section{Tensor Origin of the Relaxation Scale}
% ==========================================================

The tensor sector provides the most direct physical interpretation of the
relaxation scale. Tensor perturbations obey
\begin{equation}
  \ddot{h}_{k}
  +
  (3H + \Gamma)\, \dot{h}_{k}
  +
  \frac{k^{2}}{a^{2}} h_{k}
  =
  0.
  \label{eq:tensorrel}
\end{equation}
Using
\begin{equation}
  \Gamma(a)
  =
  \alpha(a)\, H(a),
\end{equation}
the evolution equation becomes
\begin{equation}
  \ddot{h}_{k}
  +
  (3 + \alpha) H\, \dot{h}_{k}
  +
  \frac{k^{2}}{a^{2}} h_{k}
  =
  0.
\end{equation}

Tensor modes crossing the horizon during the relaxation epoch experience a
cumulative damping given by
\begin{equation}
  T_{\mathrm{rel}}(k)
  =
  \exp
  \left[
    -\frac{1}{2}
    \int \alpha(a)\, d\ln a
  \right].
\end{equation}
For sufficiently large wavenumbers,
\begin{equation}
  k > k_{\mathrm{rel}},
\end{equation}
one obtains the approximate transfer function
\begin{equation}
  T_{\mathrm{rel}}(k)
  \simeq
  \left(
    \frac{k}{k_{\mathrm{rel}}}
  \right)^{-\alpha_{\mathrm{eff}}}.
  \label{eq:Trel}
\end{equation}

The appearance of the same characteristic scale in both tensor damping and
large-scale-structure observables strongly suggests a common geometric origin.
Consequently, Companion XII interprets Eq.~(\ref{eq:krelH}) as the fundamental
definition of the relaxation scale.
% ==========================================================
\section{Filter Interpretation of the Relaxation Scale}
% ==========================================================

The horizon definition provides a geometric interpretation of the relaxation
scale. A complementary perspective may be obtained from the dynamics of
perturbations themselves. In this picture the relaxation sector acts as a
filter that damps sufficiently short-wavelength modes.

A characteristic filter scale may therefore be defined by
\begin{equation}
  k_{\mathrm{rel}}^{(F)}
  =
  \frac{\Gamma(a)}{c_{\mathrm{eff}}(a)},
  \label{eq:krelfilter}
\end{equation}
where \(c_{\mathrm{eff}}\) denotes the effective propagation speed of
relaxation disturbances in the coupled visible--mirror system.

The precise microscopic form of \(c_{\mathrm{eff}}\) is not currently known
within the Phase-1 framework. In general one expects
\begin{equation}
  c_{\mathrm{eff}}
  =
  f
  \big(
    c_{s,+},\,
    c_{s,-},\,
    \rho_{(+)},\,
    \rho_{(-)}
  \big),
\end{equation}
where
\begin{itemize}
  \item \(c_{s,+}\) is the visible-sector sound speed,
  \item \(c_{s,-}\) is the mirror-sector sound speed,
  \item \(\rho_{(+)}\) and \(\rho_{(-)}\) are the sectoral energy densities.
\end{itemize}
The explicit construction of this effective velocity remains an open problem
and is deferred to future work.

Nevertheless, Eq.~(\ref{eq:krelfilter}) provides useful physical intuition:
relaxation becomes effective once the relaxation rate exceeds the propagation
rate of perturbations.

% ==========================================================
\section{Consistency Relation Between the Two Pictures}
% ==========================================================

At freeze-out,
\begin{equation}
  \Gamma(a_{\mathrm{fr}})
  =
  H(a_{\mathrm{fr}}).
\end{equation}
Substituting this relation into Eq.~(\ref{eq:krelfilter}) yields
\begin{equation}
  k_{\mathrm{rel}}^{(F)}
  =
  \frac{H(a_{\mathrm{fr}})}{c_{\mathrm{eff}}(a_{\mathrm{fr}})}.
\end{equation}
Expressed in comoving variables,
\begin{equation}
  k_{\mathrm{rel}}^{(F)}
  =
  \frac{a_{\mathrm{fr}} H(a_{\mathrm{fr}})}{c_{\mathrm{eff}}(a_{\mathrm{fr}})}.
\end{equation}
Using the horizon definition,
\begin{equation}
  k_{\mathrm{rel}}^{(H)}
  =
  a_{\mathrm{fr}} H(a_{\mathrm{fr}}),
\end{equation}
one finds
\begin{equation}
  k_{\mathrm{rel}}^{(F)}
  =
  \frac{k_{\mathrm{rel}}^{(H)}}{c_{\mathrm{eff}}(a_{\mathrm{fr}})}.
  \label{eq:consistency}
\end{equation}

The two descriptions therefore differ only through an order-unity factor
associated with the effective signal speed of the coupled sectors.
Companion~XII interprets this relation as evidence that both pictures describe
the same underlying physical scale. The horizon picture provides a geometric
definition, whereas the filter picture provides a dynamical interpretation.

% ==========================================================
\section{From Freeze-Out to Structure Formation}
% ==========================================================

The relaxation scale immediately propagates into the observable structure of
the Universe. Modes satisfying
\begin{equation}
  k < k_{\mathrm{rel}}
\end{equation}
remain largely unaffected by relaxation effects. Modes satisfying
\begin{equation}
  k > k_{\mathrm{rel}}
\end{equation}
experience progressive suppression.

This behavior generates the characteristic RDCC modification of growth
observables. A generic quantity
\begin{equation}
  X \in \{D, f, \gamma\}
\end{equation}
may be written schematically as
\begin{equation}
  X_{\mathrm{RDCC}}(k,a)
  =
  X(a)
  \left[
    1
    -
    \chi
    \left(
      \frac{k}{k_{\mathrm{rel}}}
    \right)^{n}
  \right],
\end{equation}
where
\begin{itemize}
  \item \(D\) denotes the growth factor,
  \item \(f\) denotes the growth rate,
  \item \(\gamma\) denotes the growth index,
  \item \(n\) characterizes the suppression shape.
\end{itemize}
The precise value of \(n\) depends on the detailed perturbative treatment and
therefore belongs to the modernization of Companion~VII.

The important result established here is that all of these observables inherit
the same relaxation scale defined by the freeze-out geometry.

% ==========================================================
\section{Connection to Companion VII}
% ==========================================================

The Boltzmann hierarchy of RDCC uses the mirror-sector contrast
\(\delta_{(-)}\) instead of a conventional cold-dark-matter perturbation. The
relaxation scale derived in the present Companion provides the missing
geometric input required for a consistent treatment of
\begin{itemize}
  \item growth suppression,
  \item velocity divergence,
  \item redshift-space distortions,
  \item scale-dependent clustering.
\end{itemize}
Companion~XII therefore supplies the theoretical foundation required for the
next-stage modernization of Companion~VII.
% ==========================================================
\section{Connection to Gravitational Slip}
% ==========================================================

The gravitational slip parameter
\begin{equation}
  \eta
  =
  \frac{\Phi}{\Psi}
\end{equation}
provides a sensitive probe of modifications to the metric sector. Within the
RDCC framework, the relaxation scale defines the characteristic transition
between scales that retain their standard metric response and scales affected
by relaxation dynamics. Consequently, the same relaxation scale appears
naturally in the gravitational slip sector.

Modes satisfying
\begin{equation}
  k > k_{\mathrm{rel}}
\end{equation}
are expected to develop the strongest deviation from the standard relation
\begin{equation}
  \Phi = \Psi.
\end{equation}
The detailed derivation of the scale-dependent slip belongs to the
modernization of Companion~XLIV. The present Companion establishes only the
geometric origin of the characteristic transition scale.

% ==========================================================
\section{Weak Lensing and Survey Observables}
% ==========================================================

Weak gravitational lensing depends on the combined lensing potential
\begin{equation}
  \Phi_{\mathrm{lens}}
  =
  \frac{\Phi + \Psi}{2}.
\end{equation}
Because both metric potentials inherit the relaxation scale, lensing
observables also become sensitive to
\begin{equation}
  k_{\mathrm{rel}}.
\end{equation}

The same scale is expected to appear in
\begin{itemize}
  \item cosmic shear,
  \item lensing convergence spectra,
  \item weak-lensing tomography,
  \item Euclid-scale survey analyses.
\end{itemize}
The relaxation scale therefore acts as a common geometric thread connecting
multiple observational sectors.

% ==========================================================
\section{The Unified Role of the Relaxation Scale}
% ==========================================================

The analysis of this Companion reveals an important structural result. The
quantity
\begin{equation}
  k_{\mathrm{rel}}
  =
  a_{\mathrm{fr}} H(a_{\mathrm{fr}})
\end{equation}
appears simultaneously in
\begin{itemize}
  \item tensor damping,
  \item stochastic gravitational-wave spectra,
  \item growth suppression,
  \item growth indices,
  \item redshift-space distortions,
  \item weak lensing,
  \item gravitational slip.
\end{itemize}
Rather than introducing separate scales for each of these phenomena, RDCC
identifies a single common origin in the freeze-out geometry of the
relaxation sector. The relaxation scale therefore represents a geometric
transition scale rather than an arbitrary fit parameter.

% ==========================================================
\section{Open Problems}
% ==========================================================

Although the origin of the relaxation scale has been clarified significantly,
several questions remain open.
\begin{enumerate}[label=\arabic*.]
  \item A complete derivation of
  \begin{equation}
    c_{\mathrm{eff}}
    =
    f\big(c_{s,+},\, c_{s,-},\, \rho_{(+)},\, \rho_{(-)}\big)
  \end{equation}
  remains unavailable.

  \item The numerical evaluation of
  \begin{equation}
    k_{\mathrm{rel}}
    =
    a_{\mathrm{fr}} H(a_{\mathrm{fr}})
  \end{equation}
  within the Phase-1 integrated background cosmology has not yet been carried
  out.

  \item The exact relationship between the horizon and filter interpretations
  of the relaxation scale remains to be fully established.

  \item The detailed propagation of the relaxation scale into the Boltzmann
  hierarchy will be developed in the modernized version of Companion~VII.
\end{enumerate}
These open questions do not invalidate the geometric interpretation developed
here, but instead define the next stage of the RDCC research program.

% ==========================================================
\section{Conclusion}
% ==========================================================

This Companion establishes the theoretical origin of the RDCC relaxation
scale. Rather than treating \(k_{\mathrm{rel}}\) as an independent
phenomenological parameter, the scale is reinterpreted as a consequence of the
freeze-out geometry of the relaxation sector.

The central result is
\begin{equation}
  k_{\mathrm{rel}}
  =
  a_{\mathrm{fr}} H(a_{\mathrm{fr}}),
\end{equation}
where
\begin{equation}
  a_{\mathrm{fr}} \simeq 10^{-4}
\end{equation}
marks the transition from relaxation-dominated evolution to the
post-relaxation regime.

The relaxation scale therefore provides a common bridge between the tensor
sector, large-scale-structure formation, gravitational lensing, redshift-space
distortions, and future survey observables. Together with Companion~II,
Companion~VI, and Companion~XI, the present Companion closes one of the
principal conceptual gaps of the RDCC framework and establishes the foundation
required for the modernization of Companion~VII and Companion~XLIV.

\end{document}